\documentclass[../../main.tex]{subfiles}
\begin{document}
% SOURCE p=214; printed=199; object=ch06-distribution-exercises-title; region=middle
\Needspace{12\baselineskip}
\subsection*{习题}
% SOURCE p=214; printed=199; object=ch06-ex35; number=35; region=middle
\begin{exercise}\label{exercise:ch06-weak-lp-topology}
对任意可测 $f,g$，有 $[cf]_p=|c|[f]_p$ 以及
\begin{align*}
[f+g]_p\leqslant2\bigl([f]_p^p+[g]_p^p\bigr)^{1/p},
\end{align*}
所以弱 $L^p$ 是向量空间。而且，“球” $\{g:[g-f]_p<r\}$，其中 $r>0$、$f\in\text{弱 }L^p$，生成弱 $L^p$ 上的一个拓扑，使它成为拓扑向量空间。
\end{exercise}
% SUPPLEMENT: id=v14-ch06-ex35-solution; parent=ch06-ex35; kind=solution+motivation+hypotheses; origin=independent-derivation
\begin{solution}
\textbf{弱三角估计与显式度量（重构新增）。}沿用 $0<p<\infty$，并按几乎处处相等识别函数。令 $N(f)=[f]_p$。$c\ne0$ 时有 $\lambda_{cf}(\alpha)=\lambda_f(\alpha/|c|)$，换元取上确界便给 $N(cf)=|c|N(f)$；$c=0$ 直接成立。由水平集的二分估计，
\begin{align*}
\alpha^p\lambda_{f+g}(\alpha)
\leqslant2^p\left[(\alpha/2)^p\lambda_f(\alpha/2)
+(\alpha/2)^p\lambda_g(\alpha/2)\right]
\leqslant2^p\bigl(N(f)^p+N(g)^p\bigr).
\end{align*}
取上确界与 $p$ 次方根，得到题中不等式。它与齐次性说明有限准范数的类对加法和数乘封闭。

为严格证明所说的球确实给出向量空间拓扑，再证明一个更强的不等式。记 $\theta=p/(p+1)\in(0,1)$。若 $A=N(f)>0$、$B=N(g)>0$ 都有限，对任意 $0<t<1$，有
\begin{align*}
\lambda_{f+g}(\alpha)\leqslant\lambda_f(t\alpha)+\lambda_g((1-t)\alpha),\qquad
N(f+g)^p\leqslant A^pt^{-p}+B^p(1-t)^{-p}.
\end{align*}
取 $t=A^\theta/(A^\theta+B^\theta)$，利用 $p-p\theta=\theta$，得
\begin{align*}
N(f+g)^p\leqslant(A^\theta+B^\theta)^{p+1},
\qquad N(f+g)^\theta\leqslant N(f)^\theta+N(g)^\theta.
\end{align*}
若 $A=0$ 或 $B=0$，相应函数为零几乎处处，最后的不等式仍成立。这里 $N(f)=0$ 当且仅当 $f=0$ 几乎处处，因为正水平集的可数并等于非零集。

因此
\begin{align*}
d(f,g)=N(f-g)^\theta
\end{align*}
在上述商空间上是平移不变度量：对称性、正定性由前述性质给出，三角不等式来自刚证明的幂次三角式。原题的球恰好是这个度量半径为 $r^\theta$ 的开球，故它们生成度量拓扑。加法连续，因为
\begin{align*}
d(f+g,f_0+g_0)\leqslant d(f,f_0)+d(g,g_0).
\end{align*}
数乘联合连续也可直接核验：
\begin{align*}
d(cf,c_0f_0)
\leqslant |c|^\theta d(f,f_0)+|c-c_0|^\theta N(f_0)^\theta.
\end{align*}
当 $(c,f)\to(c_0,f_0)$ 时，$c$ 在一个有界邻域内，右边趋零。因此这是 Hausdorff 拓扑向量空间。
\end{solution}
\begin{remark}
\textbf{动机与条件说明（重构新增）。}原弱三角式已保证向量运算封闭，但为了避免未经证明便把所有准范数球当作开球，解答用一次优化阈值分割构造出真正的度量。$p>0$ 使 $\theta$ 位于 $(0,1)$ 并使各次幂单调；$p<\infty$ 保证这些数为有限指数。必须识别几乎处处相等的函数，否则 $d$ 只是伪度量且一般不 Hausdorff。共同的测度空间使水平集比较有意义；无需有限测度或局部凸性。这里证明的是向量拓扑，不把原来的 $[\cdot]_p$ 误称为满足三角不等式的范数。
\end{remark}
% END SUPPLEMENT: id=v14-ch06-ex35-solution
% SOURCE p=214; printed=199; object=ch06-ex36; number=36; region=middle
\begin{exercise}\label{exercise:ch06-weak-strong-inclusions}
若 $f\in\text{弱 }L^p$ 且 $\mu\{x:f(x)\ne0\}<\infty$，则对所有 $q<p$ 有 $f\in L^q$。另一方面，若 $f\in(\text{弱 }L^p)\cap L^\infty$，则对所有 $q>p$ 有 $f\in L^q$。
\end{exercise}
% SUPPLEMENT: id=v14-ch06-ex36-solution; parent=ch06-ex36; kind=solution+motivation+hypotheses; origin=independent-derivation
\begin{solution}
\textbf{分别修复两个端点（重构新增）。}这里的 $L^q$ 指正指数空间：第一问为 $0<q<p$，第二问先证 $p<q<\infty$，而 $q=\infty$ 已在假设中。令 $A=[f]_p$。若 $A=0$，则 $f=0$ 几乎处处，所有结论成立。

第一问记 $M=\mu\{f\ne0\}<\infty$。$M=0$ 同样给零函数；以下设 $A,M>0$。对任意 $\alpha>0$，
\begin{align*}
\lambda_f(\alpha)\leqslant\min\{M,A^p\alpha^{-p}\}.
\end{align*}
取 $a=AM^{-1/p}>0$。由层积分公式，
\begin{align*}
\|f\|_q^q
&=q\int_0^\infty\alpha^{q-1}\lambda_f(\alpha)\,\mathrm d\alpha\\
&\leqslant qM\int_0^a\alpha^{q-1}\,\mathrm d\alpha
+qA^p\int_a^\infty\alpha^{q-p-1}\,\mathrm d\alpha\\
&=Ma^q+\frac q{p-q}A^pa^{q-p}
=\frac p{p-q}A^qM^{1-q/p}<\infty.
\end{align*}
第二问令 $B=\|f\|_\infty<\infty$。$B=0$ 时已证；否则 $\lambda_f(\alpha)=0$ 对 $\alpha\geqslant B$ 成立，故
\begin{align*}
\|f\|_q^q\leqslant qA^p\int_0^B\alpha^{q-p-1}\,\mathrm d\alpha
=\frac q{q-p}A^pB^{q-p}<\infty.
\end{align*}
这证明了全部所需包含关系。
\end{solution}
\begin{remark}
\textbf{动机与条件说明（重构新增）。}有限的非零支集控制低水平处的总测度，所以允许较小指数；本质有界性消除高水平尾部，所以允许较大指数。严格不等式 $q<p$ 或 $q>p$ 恰保证对应幂积分收敛，$q>0$ 保证低端的基本幂积分有意义。不能随意加入 $q=p$：$x^{-1/p}\chi_{(0,1)}$ 有有限非零支集且属于弱 $L^p$，但不在 $L^p$；$x^{-1/p}\chi_{(1,\infty)}$ 本质有界且属于弱 $L^p$，也不在 $L^p$。全部公式允许任意测度空间，不要求 $p,q\geqslant1$。
\end{remark}
% END SUPPLEMENT: id=v14-ch06-ex36-solution
% SOURCE p=214; printed=199; object=ch06-ex37; number=37; region=middle
\begin{exercise}\label{exercise:ch06-truncation-distributions}
证明\cref{proposition:ch06-distribution-truncation}。
\end{exercise}
% SUPPLEMENT: id=v14-ch06-ex37-solution; parent=ch06-ex37; kind=solution+motivation+hypotheses; origin=independent-derivation
\begin{solution}
\textbf{逐点分解与全部水平集（重构新增）。}在 $\{|f|\leqslant A\}$ 上，$h_A=f$、$g_A=0$；在 $\{|f|>A\}$ 上，相位模为一，$|h_A|=A$、$|g_A|=|f|-A$。因此对每个 $\alpha>0$，
\begin{align*}
\{|g_A|>\alpha\}&=\{|f|>A+\alpha\},\\
\{|h_A|>\alpha\}&=\begin{cases}\{|f|>\alpha\},&0<\alpha<A,\\\varnothing,&\alpha\geqslant A.\end{cases}
\end{align*}
分别取测度，就得到 $\lambda_{g_A}(\alpha)=\lambda_f(A+\alpha)$ 及 $\lambda_{h_A}$ 的两段公式，包含端点 $\alpha=A$。原式 $f=h_A+g_A$ 在两块区域上均直接成立；$f=0$ 处使用 $\operatorname{sgn}0=0$。可测性随相位和水平集的可测性得到。
\end{solution}
\begin{remark}
\textbf{动机与条件说明（重构新增）。}截断模长而保持相位，才能把复杂函数分成有界部分与超出阈值的余量。严格水平集在阈值处不包含模恰为 $A$ 的点，所以第二段必须从 $\alpha\geqslant A$ 开始。$A>0$ 与可测性沿用命题；无需积分或测度有限性。这里给出了本题完整解答，而不仅是回指正文补证。
\end{remark}
% END SUPPLEMENT: id=v14-ch06-ex37-solution
% SOURCE p=214; printed=199; object=ch06-ex38; number=38; region=middle
\begin{exercise}\label{exercise:ch06-dyadic-distribution}
$f\in L^p$ 当且仅当 $\sum_{k=-\infty}^{\infty}2^{kp}\lambda_f(2^k)<\infty$。
\end{exercise}
% SUPPLEMENT: id=v14-ch06-ex38-solution; parent=ch06-ex38; kind=solution+motivation+hypotheses; origin=independent-derivation
\begin{solution}
\textbf{二进区间上的两侧比较（重构新增）。}沿用 $0<p<\infty$。在 $[2^k,2^{k+1})$ 上，分布函数的单调性给
\begin{align*}
\lambda_f(2^{k+1})\leqslant\lambda_f(\alpha)\leqslant\lambda_f(2^k).
\end{align*}
乘以 $p\alpha^{p-1}$ 并积分，得到
\begin{align*}
(2^p-1)2^{kp}\lambda_f(2^{k+1})
\leqslant p\int_{2^k}^{2^{k+1}}\alpha^{p-1}\lambda_f(\alpha)\,\mathrm d\alpha
\leqslant(2^p-1)2^{kp}\lambda_f(2^k).
\end{align*}
对所有整数 $k$ 求和，非负项允许无穷求和及重编号。记 $S=\sum_{k\in\mathbb Z}2^{kp}\lambda_f(2^k)$，层积分公式给出
\begin{align*}
(1-2^{-p})S\leqslant\|f\|_p^p\leqslant(2^p-1)S.
\end{align*}
两个常数均严格正且有限，因而中间量有限当且仅当 $S$ 有限。
\end{solution}
\begin{remark}
\textbf{动机与条件说明（重构新增）。}只在几何级数高度处采样，借助分布函数单调性即可控制所有中间高度。整数指标必须双向无穷，才能同时探测零端与无穷端；只取正整数会遗漏一个端点。$0<p<\infty$ 保证比较常数有效，非负性允许无需绝对收敛先重排级数。证明不要求 $f$ 预先可积，也不需要 $\mu$ 为 $\sigma$ 有限。
\end{remark}
% END SUPPLEMENT: id=v14-ch06-ex38-solution
% SOURCE p=214; printed=199; object=ch06-ex39; number=39; region=middle-lower
\begin{exercise}\label{exercise:ch06-distribution-end-limits}
若 $f\in L^p$，则
\begin{align*}
\lim_{\alpha\to0}\alpha^p\lambda_f(\alpha)
=\lim_{\alpha\to\infty}\alpha^p\lambda_f(\alpha)=0.
\end{align*}
（先设 $f$ 为简单函数。）
\end{exercise}
% SUPPLEMENT: id=v14-ch06-ex39-solution; parent=ch06-ex39; kind=solution+motivation+hypotheses; origin=independent-derivation
\begin{solution}
\textbf{简单函数与统一的控制收敛证明（重构新增）。}沿用 $0<p<\infty$。若简单函数 $s\in L^p$，它的非零支集为有限个正值水平集的并，测度有限。因此 $\alpha\downarrow0$ 时，$0\leqslant\alpha^p\lambda_s(\alpha)\leqslant\alpha^p\mu\{s\ne0\}\to0$；而当 $\alpha$ 不小于 $\|s\|_\infty$ 时，$\lambda_s(\alpha)=0$，故另一端也成立。

对一般 $f\in L^p$，直接把所求量写成积分：
\begin{align*}
\alpha^p\lambda_f(\alpha)=\int_X\alpha^p\chi_{\{|f|>\alpha\}}\,\mathrm d\mu,
\qquad 0\leqslant\alpha^p\chi_{\{|f|>\alpha\}}\leqslant|f|^p\in L^1.
\end{align*}
当 $\alpha\downarrow0$ 时，被积函数不超过 $\alpha^p$，逐点趋零；当 $\alpha\to\infty$ 时，对每个 $|f(x)|<\infty$ 的点，示性函数最终为零。$f$ 有限几乎处处，所以控制收敛定理在两端都适用。若要把序列形式写成参数极限，对任意趋向相应端点的正数序列应用同一论证即可；序列判据给出所求两个极限。
\end{solution}
\begin{remark}
\textbf{动机与条件说明（重构新增）。}Chebyshev 只给有界的 $\alpha^p\lambda_f(\alpha)$，而真正的 $L^p$ 可积性提供一个固定的可积控制函数，从而使端点贡献消失。弱 $L^p$ 中不能保证这一点，例如 $x^{-1/p}$ 的该量恒为一。$0<p<\infty$ 与 $|f|^p$ 的可积性是关键，不使用范数三角不等式，因此也包含 $p<1$。这补全了前面 Stieltjes 分部积分所需的边界说明。
\end{remark}
% END SUPPLEMENT: id=v14-ch06-ex39-solution
% SOURCE p=214; printed=199; object=ch06-ex40; number=40; region=lower
\begin{exercise}\label{exercise:ch06-decreasing-rearrangement}
\begin{definition*}[递减重排]\phantomsection\label{definition:ch06-decreasing-rearrangement}
若 $f$ 为 $X$ 上的可测函数，其\textbf{递减重排}是函数 $f^*:(0,\infty)\to[0,\infty]$，定义为
\begin{align*}
f^*(t)=\inf\{\alpha:\lambda_f(\alpha)\leqslant t\},\qquad\inf\varnothing=\infty.
\end{align*}
\end{definition*}
% SUPPLEMENT: id=v14-ch06-rearrangement-definition-motivation; parent=ch06-ex40; kind=motivation+hypotheses; origin=independent-derivation
\begin{remark}
\textbf{定义的动机与边界约定（重构新增）。}下确界中的参数沿用 $\lambda_f$ 的定义域，故为 $\alpha>0$。$f^*(t)$ 寻找能把超过该高度的测度降到 $t$ 以下的最低阈值；于是把大值排在前面，只保留各高度所占测度。若没有这样的高度，约定下确界为无穷，保留一般可测函数的情况。原题 (a) 会出现 $\lambda_f(f^*(t))$ 与 $f^*(\lambda_f(\alpha))$；内层值可能为零，超出刚写出的开区间定义域。为使这两式覆盖零值情形，明确作边界延拓
\begin{align*}
\lambda_f(0)=\mu\{|f|>0\}=\lim_{\alpha\downarrow0}\lambda_f(\alpha),
\qquad f^*(0)=\inf\{\alpha>0:\lambda_f(\alpha)=0\}.
\end{align*}
两者均允许无穷。原定义在正参数上完全保持；后面的复合积分继续使用 $\phi(0)=0$ 的约定。这里 $t$ 是在重排空间中的长度参数，而 $\alpha$ 是函数高度，二者不可互换。
\end{remark}
% END SUPPLEMENT: id=v14-ch06-rearrangement-definition-motivation
\begin{enumerate}
\item\label{item:ch06-rearrangement-inverse} $f^*$ 递减。若 $f^*(t)<\infty$，则 $\lambda_f(f^*(t))\leqslant t$；若 $\lambda_f(\alpha)<\infty$，则 $f^*(\lambda_f(\alpha))\leqslant\alpha$。
\item\label{item:ch06-rearrangement-equimeasurable} $\lambda_f=\lambda_{f^*}$，其中 $\lambda_{f^*}$ 相对于 $(0,\infty)$ 上的 Lebesgue 测度定义。
\item\label{item:ch06-rearrangement-integrals} 若对所有 $\alpha>0$，$\lambda_f(\alpha)<\infty$，并且 $\lim_{\alpha\to\infty}\lambda_f(\alpha)=0$（从而对所有 $t>0$，$f^*(t)<\infty$），而 $\phi\geqslant0$ 是 $(0,\infty)$ 上的 Borel 可测函数，则
\begin{align*}
\int_X\phi\circ|f|\,\mathrm d\mu=\int_0^\infty\phi\circ f^*(t)\,\mathrm dt.
\end{align*}
特别地，对 $0<p<\infty$，$\|f\|_p=\|f^*\|_p$。
\item\label{item:ch06-rearrangement-weak-norm} 若 $0<p<\infty$，则 $[f]_p=\sup_{t>0}t^{1/p}f^*(t)$。
\item\label{item:ch06-rearrangement-graph} “重排”这一名称来自 $f$ 为 $(0,\infty)$ 上的非负函数的情形。为看出它为何恰当，选一个在 $(0,\infty)$ 上取四个或五个不同值的阶梯函数，并画出 $f$ 与 $f^*$ 的图像。
\end{enumerate}
\end{exercise}
% SUPPLEMENT: id=v14-ch06-ex40-solution; parent=ch06-ex40; kind=solution+motivation+hypotheses+figure; origin=independent-derivation
\begin{solution}
\textbf{全部五问的证明与作图（重构新增）。}
\begin{enumerate}
\item 若 $0<t_1<t_2$，则 $\{\alpha>0:\lambda_f(\alpha)\leqslant t_1\}\subset\{\alpha>0:\lambda_f(\alpha)\leqslant t_2\}$，故其下确界满足 $f^*(t_1)\geqslant f^*(t_2)$。因此 $f^*$ 非增，作为扩展非负函数可测。

设 $b=f^*(t)<\infty$。允许集合非空；对任意 $\beta>b$，下确界性质给其中某个 $a<\beta$，再用 $\lambda_f$ 非增，得到 $\lambda_f(\beta)\leqslant\lambda_f(a)\leqslant t$。令 $\beta\downarrow b$，若 $b>0$ 用右连续性，若 $b=0$ 用前述边界延拓，都得 $\lambda_f(b)\leqslant t$。其次，令 $s=\lambda_f(\alpha)<\infty$。若 $s>0$，$\alpha$ 本身就在定义 $f^*(s)$ 的允许集合内；若 $s=0$，则在定义 $f^*(0)$ 的集合内。所以两种情况均有 $f^*(s)\leqslant\alpha$。
\item 对每个 $\alpha>0$，有精确的集合等式
\begin{align*}
\{t>0:f^*(t)>\alpha\}=\{t>0:t<\lambda_f(\alpha)\}.
\end{align*}
若 $\lambda_f(\alpha)\leqslant t$，则 $\alpha$ 已是允许阈值，故 $f^*(t)\leqslant\alpha$。反之，若 $\lambda_f(\alpha)>t$，由右连续性可取 $\beta>\alpha$ 使 $\lambda_f(\beta)>t$；当 $\lambda_f(\alpha)=\infty$ 时同样成立，因为右极限为无穷。所有不超过 $\beta$ 的正高度都不是允许阈值，所以 $f^*(t)\geqslant\beta>\alpha$。右边集合是长度为 $\lambda_f(\alpha)$ 的开区间，长度可为零或无穷，故 $\lambda_{f^*}(\alpha)=\lambda_f(\alpha)$。
\item 因 $\lambda_f(\alpha)\to0$，对每个 $t>0$ 可找到允许阈值，故 $f^*(t)<\infty$。由 (b)，$f$ 与 $f^*$ 有相同且在各正高度有限的分布函数，因此分别对它们应用\cref{proposition:ch06-distribution-pushforward}，得到
\begin{align*}
\int_X\phi(|f|)\,\mathrm d\mu
=-\int_0^\infty\phi(\alpha)\,\mathrm d\lambda_f(\alpha)
=-\int_0^\infty\phi(\alpha)\,\mathrm d\lambda_{f^*}(\alpha)
=\int_0^\infty\phi(f^*(t))\,\mathrm dt.
\end{align*}
两次都把 $\phi$ 零延拓到原点。取 $\phi(\alpha)=\alpha^p$ 并取 $p$ 次方根即得范数等式，允许共同的无穷值。
\item 令 $N=[f]_p=\sup_{\alpha>0}\alpha\lambda_f(\alpha)^{1/p}$，以及 $M=\sup_{t>0}t^{1/p}f^*(t)$。若 $0<t<\lambda_f(\alpha)$，由 (b) 的集合等式有 $f^*(t)>\alpha$，从而 $M\geqslant\alpha t^{1/p}$。令 $t\uparrow\lambda_f(\alpha)$，得 $M\geqslant\alpha\lambda_f(\alpha)^{1/p}$；当该分布值为无穷时取任意大的 $t$，当它为零时该不等式平凡。因此 $M\geqslant N$。

反之，固定 $t>0$。若 $f^*(t)>0$，对每个 $0<\alpha<f^*(t)$，(b) 给 $t<\lambda_f(\alpha)$，故 $\alpha t^{1/p}\leqslant N$。令 $\alpha\uparrow f^*(t)$，包括 $f^*(t)=\infty$ 的情况，得到 $t^{1/p}f^*(t)\leqslant N$；若 $f^*(t)=0$，同一结论平凡。取 $t$ 的上确界，得 $M\leqslant N$，故相等。
\item 在正半轴上取五值阶梯函数
\begin{align*}
f(x)=\begin{cases}1,&0<x<1,\\4,&1\leqslant x<2,\\2,&2\leqslant x<3,\\3,&3\leqslant x<4,\\0,&x\geqslant4.\end{cases}
\end{align*}
其分布函数在 $0<\alpha<1$、$1\leqslant\alpha<2$、$2\leqslant\alpha<3$、$3\leqslant\alpha<4$、$\alpha\geqslant4$ 上依次为 $4,3,2,1,0$。按下确界定义，
\begin{align*}
f^*(t)=\begin{cases}4,&0<t<1,\\3,&1\leqslant t<2,\\2,&2\leqslant t<3,\\1,&3\leqslant t<4,\\0,&t\geqslant4.\end{cases}
\end{align*}
所以各高度占有的长度不变，只按高度从大到小排列。图中的实心点取该端点值，空心点表示不包含该端点；零点本不在定义域。
\begin{figure}[H]
\centering
\begin{tikzpicture}[x=0.8cm,y=0.65cm,>=stealth]
\begin{scope}
\draw[->] (0,0)--(5,0) node[right] {$x$};
\draw[->] (0,0)--(0,4.8) node[above] {$f(x)$};
\foreach \k in {1,2,3,4}{\draw (\k,0)--(\k,-0.09) node[below] {\small $\k$};\draw (0,\k)--(-0.09,\k) node[left] {\small $\k$};}
\draw[thick,blue] (0,1)--(1,1) (1,4)--(2,4) (2,2)--(3,2) (3,3)--(4,3) (4,0)--(4.8,0);
\foreach \x/\y in {0/1,1/1,2/4,3/2,4/3}{\draw[blue,fill=white] (\x,\y) circle (2pt);}
\foreach \x/\y in {1/4,2/2,3/3,4/0}{\fill[blue] (\x,\y) circle (2pt);}
\end{scope}
\begin{scope}[xshift=6cm]
\draw[->] (0,0)--(5,0) node[right] {$t$};
\draw[->] (0,0)--(0,4.8) node[above] {$f^*(t)$};
\foreach \k in {1,2,3,4}{\draw (\k,0)--(\k,-0.09) node[below] {\small $\k$};\draw (0,\k)--(-0.09,\k) node[left] {\small $\k$};}
\draw[thick,teal] (0,4)--(1,4) (1,3)--(2,3) (2,2)--(3,2) (3,1)--(4,1) (4,0)--(4.8,0);
\foreach \x/\y in {0/4,1/4,2/3,3/2,4/1}{\draw[teal,fill=white] (\x,\y) circle (2pt);}
\foreach \x/\y in {1/3,2/2,3/1,4/0}{\fill[teal] (\x,\y) circle (2pt);}
\end{scope}
\end{tikzpicture}
\caption{习题中的阶梯函数及其递减重排（重构新增示例与作图）}
\label{figure:ch06-rearrangement-step-example}
\end{figure}
\end{enumerate}
\end{solution}
\begin{remark}
\textbf{动机与逐条件说明（重构新增）。}单调性让重排成为分布函数的广义逆；右连续性则使严格水平集恰好对应开区间，若省略它，阈值等号处的推理便可能失效。无穷下确界约定保证 (a)、(b)、(d) 覆盖无有限阈值的情形；零参数延拓则处理零函数和有界函数可能出现的边界值。(c) 的正水平有限保证 Stieltjes 测度局部有限，尾分布趋零保证重排逐点有限，非负 Borel $\phi$ 允许单调积分并通过零延拓忽略原函数的零值区域。范数等式用 $0<p<\infty$ 的层高度函数；(d) 也要求正有限指数，但不要求分布函数有限或尾部消失，证明已逐一处理无穷值。作图选非负函数使图像高度直接代表模值；对复函数重排依赖的是 $|f|$，不保留相位。
\end{remark}
% END SUPPLEMENT: id=v14-ch06-ex40-solution
\end{document}
